% Source: https://artofproblemsolving.com/community/c3h2924123_geo_proof_problems
% Stable identity: aops-topic-2924123-q01
\emph{Case 1: $O$ lies on $CA$.} Triangle $OBC$ is isosceles ($OB=OC$), so $\angle OCB=\angle OBC$. The exterior angle at $O$ is $\angle AOB=\angle OCB+\angle OBC=2\angle ACB$.

\emph{General case.} Draw the diameter $CD$.
\begin{asy}
usepackage("amsmath");
size(3cm);
defaultpen(fontsize(10pt));
pair O=(0,0),A=dir(250),B=dir(-20),C=dir(110);pair A2=-C;draw(circle(O,1));draw(C--A--O--B--C);draw(C--A2,dashed);dot(O);
label("$O$",O,E);label("$A$",A,SW);label("$B$",B,E);label("$C$",C,N);label("$D$",A2,S);
\end{asy}

If $O$ is inside $\angle ACB$, then $\angle ACB=\angle ACD+\angle DCB=\dfrac12\angle AOD+\dfrac12\angle DOB=\dfrac12\angle AOB$ by Case 1 applied twice. If $O$ is outside, subtract instead of adding. $\blacksquare$
